% Intended LaTeX compiler: pdflatex
\documentclass[a4paper,12pt]{article}
\usepackage[utf8]{inputenc}
\usepackage[T1]{fontenc}
\usepackage{graphicx}
\usepackage{longtable}
\usepackage{wrapfig}
\usepackage{rotating}
\usepackage[normalem]{ulem}
\usepackage{amsmath}
\usepackage{amssymb}
\usepackage{capt-of}
\usepackage{hyperref}
\usepackage{\string~/.emacs.d/common}
\author{mahmood}
\date{\textit{<2022-06-18 Sat 14:15:00>}}
\title{sagemath}
\hypersetup{
 pdfauthor={mahmood},
 pdftitle={sagemath},
 pdfkeywords={},
 pdfsubject={sagemath},
 pdfcreator={Emacs 30.0.50 (Org mode 9.6)}, 
 pdflang={English}}
\begin{document}

\maketitle
\begin{note}
video demonstrating sagemath+org-mode workflow:\\[0pt]
\url{https://www.youtube.com/watch?v=cd4\_O3TqArs}\\[0pt]
\end{note}
sagemath is an open source alternative to Mathematica written in python, it preparses python code to allow for extra syntactic sugars\\[0pt]
its basically a python interface to a collection of mathematical software like Maxima/sympy\\[0pt]
\section{plotting}
\label{sec:orgd2dd39a}
\subsection{area below functions}
\label{sec:orgcadc4fb}
\begin{lstlisting}[language=sage,label= ,caption= ,captionpos=b,numbers=none]
interval = [1, 8.5]
interval_fill = [2, 8]
f(x) = (x-3)*(x-5)*(x-7)+40
import os.path

try: _ = open_left
except: open_left = False
try: _ = open_right
except: open_right = False

if not open_left:
  my_plot = line([(interval_fill[0],0),(interval_fill[0],f(interval_fill[0]))], color='black')
if not open_right:
  my_plot += line([(interval_fill[1], 0),(interval_fill[1], f(interval_fill[1]))], color='black')
my_plot += plot(f, (interval_fill[0], interval_fill[1]), fill=True) + plot(f, (interval[0], interval[1]), thickness=2)
my_plot.save(filename=os.path.expanduser(filename), transparent=True)
print(os.path.expanduser(filename))

open_left = False
open_right = False
\end{lstlisting}

\includesvg{/Users/mahmooz/brain/out/Mqv2WA}\\[0pt]

\begin{lstlisting}[language=sage,label= ,caption= ,captionpos=b,numbers=none]
interval = [2, 8.5]
interval_fill = [3, 8]
f(x) = 1/x^2
import os.path

try: _ = open_left
except: open_left = False
try: _ = open_right
except: open_right = False

if not open_left:
  my_plot = line([(interval_fill[0],0),(interval_fill[0],f(interval_fill[0]))], color='black')
if not open_right:
  my_plot += line([(interval_fill[1], 0),(interval_fill[1], f(interval_fill[1]))], color='black')
my_plot += plot(f, (interval_fill[0], interval_fill[1]), fill=True) + plot(f, (interval[0], interval[1]), thickness=2)
my_plot.save(filename=os.path.expanduser(filename), transparent=True)
print(os.path.expanduser(filename))

open_left = False
open_right = False
\end{lstlisting}

\includesvg{/Users/mahmooz/brain/out/c4XFj9}\\[0pt]

\subsection{solid of revolution}
\label{sec:org599c331}
consider the solid of revolution of the function \(e^x\) rotated around the \(y\) axis\\[0pt]

\begin{lstlisting}[language=sage,label= ,caption= ,captionpos=b,numbers=none]
my_plot = revolution_plot3d(e^x, (x,0,1), show_curve=True, parallel_axis='z', figsize=5)

\end{lstlisting}

\subsection{parametric 3d plot}
\label{sec:org737b2e8}
\begin{lstlisting}[language=sage,label= ,caption= ,captionpos=b,numbers=none]
my_plot = parametric_plot3d([cos(x),sin(x),x/10], (x,0,4*pi), color='red', viewer='threejs')

\end{lstlisting}

\subsection{more}
\label{sec:orgd47cf8f}
\begin{lstlisting}[language=sage,label= ,caption= ,captionpos=b,numbers=none]
x,y,z = var('x,y,z')
f(x,y,z)=x^2+y^2+z^2/4
my_plot = Graphics()
for k in [1,2,3,4]:
    my_plot += implicit_plot3d(f(x,y,z)==k,(x,-2,2),(y,-2,2),(z,-5,5),opacity=0.4)

\end{lstlisting}

\subsection{animation}
\label{sec:org3c0c8df}
\begin{lstlisting}[language=sage,label= ,caption= ,captionpos=b,numbers=none]
def build_frame(t):
    e = parametric_plot3d([sin(x-t), 0, x], (x, 0, 2*pi), color='red')
    m = parametric_plot3d([0, -sin(x-t), x], (x, 0, 2*pi), color='green')
    return e + m
frames = [build_frame(t) for t in (0, pi/32, pi/16, .., 2*pi)]
my_plot = animate(frames).interactive()

\end{lstlisting}

\begin{lstlisting}[language=sage,label= ,caption= ,captionpos=b,numbers=none]
sines = [plot(c*sin(x), (-2*pi,2*pi), color=Color(c,0,0), ymin=-1, ymax=1,transparent=False) for c in sxrange(0,1,.1)]
a = animate(sines)
a.save(os.path.expanduser(f))
print(f)
\end{lstlisting}

\url{file:///Users/mahmooz/brain/out/IRCODD.gif}\\[0pt]

\begin{lstlisting}[language=sage,label= ,caption= ,captionpos=b,numbers=none]
sin_frames = []
circ_frames = []
circ_x = -1.5 # offset center of the circle
circ_y = 0
for i in srange(0,2*pi,0.2):
    singraph = point((i,sin(i)), color="green", size=50)
    singraph += plot(sin(x),(0,2*pi), xmin=0, xmax=7, ymin=-1, ymax=1, figsize=[2,2], axes=False)
    unitcircle = point((cos(i)+circ_x,sin(i)+circ_y), color="green", size=50)
    unitcircle += circle((circ_x,circ_y),1, color="blue", figsize=[2,2], axes=False)
    sin_frames.append(singraph)
    circ_frames.append(unitcircle)
A1 = animate(sin_frames)
A2 = animate(circ_frames)
(A1 + A2).save(os.path.expanduser(f))
print(f)
\end{lstlisting}

\url{file:///Users/mahmooz/brain/out/Sg63BY.gif}\\[0pt]

\begin{lstlisting}[language=sage,label= ,caption= ,captionpos=b,numbers=none]
f(x) = 5*x - x^2/8 + x^3/1200
df(x) = diff(f(x),x)
max = 100
p = plot(f(x),x,0,max)
lblp = text("$y = " + latex(f(x)) + "$",[40,100],fontsize=14)
lbldp = text("$y' = " + latex(df(x)) + "$",[40,80],fontsize=14,rgbcolor='#006000')
ga = []
for argx in srange(0,max,5):
  dp = plot(f(argx) + (df(argx)*(x-argx)),x,0,max,color="#006000")
  xp = point((argx,f(argx)),rgbcolor='#800000')    
  ga.append(p+lblp+lbldp+dp+xp)
a = animate(ga,ymin=0,ymax=max,axes_labels=['x','y'],fontsize=12,figsize=(4,3))
a.save(os.path.expanduser(out_file))
print(out_file)
\end{lstlisting}

\url{file:///Users/mahmooz/brain/out/yAyHE1.gif}\\[0pt]

\subsection{area between functions}
\label{sec:orgfbdc8ca}
\begin{lstlisting}[language=sage,label= ,caption= ,captionpos=b,numbers=none]
f(x) = -1/8*x^3+x
g(x) = 1/2*x
h(x) = 1/4*x^2
my_plot = plot([f,g,h],-2,3,fill={0:g,1:h,2:0})

\end{lstlisting}

\includesvg{/Users/mahmooz/brain/out/9s9RNl}\\[0pt]

\subsection{more filling}
\label{sec:org5eda63f}
\begin{lstlisting}[language=sage,label= ,caption= ,captionpos=b,numbers=none]
my_plot = plot(1.13*log(x), 1, 100, fill = lambda a: nth_prime(a)/floor(a), fillcolor = 'red')

\end{lstlisting}

\includesvg{/Users/mahmooz/brain/out/efTJka}\\[0pt]

\begin{lstlisting}[language=sage,label= ,caption= ,captionpos=b,numbers=none]
#plot(x,(x,0,1),fill=x^2)
my_plot = plot([x,x^2],(x,-0.5,1),fill={0:[1]})

\end{lstlisting}

\includesvg{/Users/mahmooz/brain/out/gYDdhB}\\[0pt]

\section{examples}
\label{sec:org84239d7}
\subsection{recursive sum with sage}
\label{sec:org66ce95f}
this also shows a limitation of sagemath's, i thought something as simple as a quick 1-liner\\[0pt]
using the sum function would do the job\\[0pt]
\begin{lstlisting}[language=sage,label= ,caption= ,captionpos=b,numbers=none]
def ev(self, x):
  if x == 0:
    return 1
  else:
    return 1 + sum(self(j) for j in range(x))
f = function("f", nargs=1, eval_func=ev)
[f(i) for i in range(12)]
\end{lstlisting}

[1, 2, 4, 8, 16, 32, 64, 128, 256, 512, 1024, 2048]\\[0pt]
\subsection{recursive sum but pure python}
\label{sec:org1db9681}
\begin{lstlisting}[language=Python,label= ,caption= ,captionpos=b,numbers=none]
def f(n):
  if n == 0:
    return 1
  else:
    total_sum = 0
    for i in range(n):
      total_sum += f(i)
    return total_sum + 1
print([f(i) for i in range(12)])
\end{lstlisting}

\begin{verbatim}
[1, 2, 4, 8, 16, 32, 64, 128, 256, 512, 1024, 2048]
\end{verbatim}
\subsection{vectors of variables}
\label{sec:org5e5af58}
\begin{lstlisting}[language=sage,label= ,caption= ,captionpos=b,numbers=none]
b = list(var("b", n=8, latex_name='\\overline{x}'))
b[0]=123
show(b)
\end{lstlisting}

[123, b1, b2, b3, b4, b5, b6, b7]\\[0pt]
\subsection{indexed variables}
\label{sec:org653621e}
consider the following summation:\\[0pt]

\[ \sum_{i=1}^{k} a_i \cdot d = \dfrac{k}{2}[2a_1 + (k - 1)d] \]\\[0pt]

here \(a_i\) is always dependant on the index \(i\) and we cant define a variable number of variables\\[0pt]
so we define \(a\) as a function of \(i\)\\[0pt]

\begin{lstlisting}[language=sage,label= ,caption= ,captionpos=b,numbers=none]
a = function("a")
d = var("d")
j, i, k = var("j, i, k", domain="integer")
S = sum(a(j)*d, j, 1, k+1)
S(j=10)
\end{lstlisting}

d*sum(a(10), 10, 1, k + 1)\\[0pt]
\subsection{symbolic functions}
\label{sec:org71f0d76}
just like you can define a symbolic variable, you can also define yet-unknown functions!\\[0pt]
\begin{lstlisting}[language=sage,label= ,caption= ,captionpos=b,numbers=none]
T = var('T')
a = function('a')
R = var('R')
v = var('v')
b = var('b')
p = R*T/v - a(T)/v/(v+b)
diff(p, T)
\end{lstlisting}

R/v - diff(a(T), T)/((b + v)*v)\\[0pt]
\subsection{boolean functions}
\label{sec:orgfc3ed17}
\begin{lstlisting}[language=sage,label= ,caption= ,captionpos=b,numbers=none]
from sage.crypto.boolean_function import BooleanFunction
f = BooleanFunction([0,0,0,1,0,1,1,1])
f.truth_table()
f.linear_structures().list()
\end{lstlisting}

\begin{lstlisting}[language=sage,label= ,caption= ,captionpos=b,numbers=none]
from sage.crypto.boolean_function import BooleanFunction
B.<a, y> = BooleanPolynomialRing()
f = BooleanFunction(-a*-y)
f2 = BooleanFunction([1,0,0,0])
f.algebraic_normal_form()
f2.truth_table()
f.truth_table()
[[f(i), f2(i)] for i in range(4)]
\end{lstlisting}
\section{limitations}
\label{sec:org181707a}

consider the following equation:\\[0pt]

\[ t = x^2 + 5 \]\\[0pt]

we use this kind of equation when integrating by substituion\\[0pt]

unfortunately sagemath cant derivate this with respect to \(t\) as the following snippet gives an error\\[0pt]

\begin{lstlisting}[language=sage,label= ,caption= ,captionpos=b,numbers=none]
t = x^2 + 5
t.derivative(t)
\end{lstlisting}

but weirdly enough it can handle the inverse integral\ldots{}\\[0pt]

\begin{lstlisting}[language=sage,label= ,caption= ,captionpos=b,numbers=none]
t = x^2 + 5
t.integrate(t)
\end{lstlisting}

1/2*(x\textsuperscript{2} + 5)\textsuperscript{2}\\[0pt]

\subsection{wrong limit}
\label{sec:org2ac9eaf}

consider the limit:\\[0pt]
\[
  \lim_{x\to\infty} \frac{1}{\tan^2 x + 1}
\]\\[0pt]
which is indeterminate, but sage apparently doesnt know that cuz it gives a result:\\[0pt]

\begin{lstlisting}[language=sage,label= ,caption= ,captionpos=b,numbers=none]
f = 1/(1+tan(x)^2)
f.limit(x=oo)
\end{lstlisting}
which is obviously incorrect..\\[0pt]

\subsection{i found this problem on a discord server dedicated to maths}
\label{sec:org8386708}

\[ \sum_{i=0}^{10}\left[\prod_{j=i}^{1}(10-j)\right] = 108 \]\\[0pt]

i tried replicating it in sagemath but unfortunately sagemath cant handle this either\\[0pt]

\begin{lstlisting}[language=sage,label= ,caption= ,captionpos=b,numbers=none]
j,i = var('j,i')
from sage.calculus.calculus import symbolic_product
sum(symbolic_product(10-j, j, i, 1), i, 0, 10)
\end{lstlisting}

this block of code makes perfect sense to me but idk why it doesnt work\\[0pt]

\subsection{limit of sequence}
\label{sec:org8445276}

consider the following limit of sequence:\\[0pt]
\[
  \lim_{n\to\infty} \frac{3n}{\sqrt{2n^4-n^2}} + \frac{3n}{\sqrt{2n^4-4n^2}} + \frac{3n}{\sqrt{2n^4-9n^2}} + \cdots + \frac{3n}{\sqrt{2n^4-n^4}} 
\]\\[0pt]
i tried implementing this in sagemath but it cant solve it properly:\\[0pt]

\begin{lstlisting}[language=sage,label= ,caption= ,captionpos=b,numbers=none]
k,n = var('k,n')
s = sum((3*n)/sqrt(2*n^4-k^2*n^2), k, 1, n)
s.limit(n=Infinity)
\end{lstlisting}

3*limit(n*sum(1/sqrt(-k\textsuperscript{2}*n\textsuperscript{2} + 2*n\textsuperscript{4}), k, 1, n), n, +Infinity)\\[0pt]

\begin{note}
i asked how to deal with this on \href{https://ask.sagemath.org/question/62497/limit-of-sequence/}{sagemath's forums}\\[0pt]
\end{note}

however, it was able to handle a simpler sequence:\\[0pt]

\[
  \lim_{n\to\infty} \frac{n}{2}\left(\frac{3}{(n+1)^2} + \frac{3}{(n+2)^2} + \cdots + \frac{3}{(n+n)^2}\right)
\]\\[0pt]

\begin{lstlisting}[language=sage,label= ,caption= ,captionpos=b,numbers=none]
%display plain
k,n = var('k,n')
s = n/2 * sum(3/(n+k)^2, k, 1, n)
s.limit(n=Infinity)
%display unicode_art
k,n = var('k,n')
s = n/2 * sum(3/(n+k)^2, k, 1, n)
s.limit(n=Infinity)
\end{lstlisting}

3/2*limit(n*(harmonic\textsubscript{number}(2*n, 2) - harmonic\textsubscript{number}(n, 2)), n, +Infinity)\\[0pt]
3/4\\[0pt]

weirdly tho we only get the correct answer when using \texttt{unicode\_art}\\[0pt]

we try sympy:\\[0pt]

\begin{lstlisting}[language=Python,label= ,caption= ,captionpos=b,numbers=none]
from sympy import limit_seq, Sum
from sympy.abc import n, k
limit_seq(n/2 * Sum(3/(n+k)**2, (k, 1, n)), n)
\end{lstlisting}

this takes forever to run, you can interrupt and rerun the last line and it would output some weird result\\[0pt]

im really thinking about switching to mathematica.\\[0pt]

\section{discrete math}
\label{sec:org587ee9f}
\subsection{defining sets}
\label{sec:org7de0087}

consider the following set:\\[0pt]

\[ A = \{x \in \mathbb{R} \mid -10 \leq x \leq 10\} \]\\[0pt]

this set is infinite as there is an infinite amount of real numbers in the interval \([-10, 10]\), to define such a set in sagemath we use \texttt{RealSet}:\\[0pt]

\begin{lstlisting}[language=sage,label= ,caption= ,captionpos=b,numbers=none]
A = RealSet.closed(-10, 10)
A.contains(5)
A.contains(-5)
A.contains(10)
A.contains(20)
\end{lstlisting}

im not sure how to do the same with other number fields like integers or rationals\\[0pt]
\section{linear algebra}
\label{sec:org69d7f10}
for reference: \href{./sage\_linear\_algebra\_refcard.pdf}{refcard}\\[0pt]
\subsection{matrices}
\label{sec:orge024463}
\subsubsection{guassian elimination}
\label{sec:orga3f2f3b}
this method reduces matrices with guassian elimination method step by step\\[0pt]
\begin{lstlisting}[language=sage,label=lst:orgd289af0,caption= ,captionpos=b,numbers=none]
def gauss_method(M,rescale_leading_entry=False):
    """Describe the reduction to echelon form of the given matrix of rationals.

    M  matrix of rationals   e.g., M = matrix(QQ, [[..], [..], ..])
    rescale_leading_entry=False  boolean  make the leading entries to 1's

    Returns: None.  Side effect: M is reduced.  Note: this is echelon form, 
    not reduced echelon form; this routine does not end the same way as does 
    M.echelon_form().

    """
    num_rows=M.nrows()
    num_cols=M.ncols()
    ld(M)

    col = 0   # all cols before this are already done
    for row in range(0,num_rows): 
        # ?Need to swap in a nonzero entry from below
        while (col < num_cols
               and M[row][col] == 0): 
            for i in M.nonzero_positions_in_column(col):
                if i > row:
                    print(" swap row",row+1,"with row",i+1)
                    M.swap_rows(row,i)
                    ld(M)
                    break     
            else:
                col += 1

        if col >= num_cols:
            break

        change_flag = False
        # Now guaranteed M[row][col] != 0
        if (rescale_leading_entry
            and M[row][col] != 1):
            fm(r"take {1/M[row][col]} times row {row+1}")
            # print(" take",1/M[row][col],"times row",row+1)
            M.rescale_row(row,1/M[row][col])
            ld(M)
            change_flag=False
        for changed_row in range(row+1,num_rows):
            if M[changed_row][col] != 0:
                change_flag=True
                factor=-1*M[changed_row][col]/M[row][col]
                print(fr'take ${latex(factor)}$ times row {row+1} plus row {changed_row+1}')
                M.add_multiple_of_row(changed_row,row,factor)
        if change_flag:
            ld(M)
            col +=1
\end{lstlisting}

example usage:\\[0pt]
\begin{lstlisting}[language=sage,label= ,caption= ,captionpos=b,numbers=none]

print("first matrix:")
R.<a> = PolynomialRing(QQ)
MS = MatrixSpace(SR, 3, 4)
m = MS.matrix([[1,a,-1,2],
               [2,-1,a,5],
               [1,10,-6,1]])
gauss_method(m)
print("second matrix:")
m = matrix(QQ, 4, [
    3, 0, 3, -2,
    -3, 3, 0, 2,
    0, -2, 2, 0,
    0, 3, -3, 0
])
gauss_method(m)
\end{lstlisting}

first matrix:\\[0pt]
\[\left[\begin{array}{rrrr} 1 & a & -1 & 2 \\ 2 & -1 & a & 5 \\ 1 & 10 & -6 & 1 \end{array}\right]\]\\[0pt]
take \(-2\) times row 1 plus row 2\\[0pt]
take \(-1\) times row 1 plus row 3\\[0pt]
\[\left[\begin{array}{rrrr} 1 & a & -1 & 2 \\ 0 & -2 \, a - 1 & a + 2 & 1 \\ 0 & -a + 10 & -5 & -1 \end{array}\right]\]\\[0pt]
take \(-\frac{a - 10}{2 \, a + 1}\) times row 2 plus row 3\\[0pt]
\[\left[\begin{array}{rrrr} 1 & a & -1 & 2 \\ 0 & -2 \, a - 1 & a + 2 & 1 \\ 0 & 0 & -\frac{{\left(a + 2\right)} {\left(a - 10\right)}}{2 \, a + 1} - 5 & -\frac{a - 10}{2 \, a + 1} - 1 \end{array}\right]\]\\[0pt]
second matrix:\\[0pt]
\[\left[\begin{array}{rrrr} 3 & 0 & 3 & -2 \\ -3 & 3 & 0 & 2 \\ 0 & -2 & 2 & 0 \\ 0 & 3 & -3 & 0 \end{array}\right]\]\\[0pt]
take \(1\) times row 1 plus row 2\\[0pt]
\[\left[\begin{array}{rrrr} 3 & 0 & 3 & -2 \\ 0 & 3 & 3 & 0 \\ 0 & -2 & 2 & 0 \\ 0 & 3 & -3 & 0 \end{array}\right]\]\\[0pt]
take \(\frac{2}{3}\) times row 2 plus row 3\\[0pt]
take \(-1\) times row 2 plus row 4\\[0pt]
\[\left[\begin{array}{rrrr} 3 & 0 & 3 & -2 \\ 0 & 3 & 3 & 0 \\ 0 & 0 & 4 & 0 \\ 0 & 0 & -6 & 0 \end{array}\right]\]\\[0pt]
take \(\frac{3}{2}\) times row 3 plus row 4\\[0pt]
\[\left[\begin{array}{rrrr} 3 & 0 & 3 & -2 \\ 0 & 3 & 3 & 0 \\ 0 & 0 & 4 & 0 \\ 0 & 0 & 0 & 0 \end{array}\right]\]\\[0pt]
\subsubsection{matrices with variables}
\label{sec:org29c164e}
to be able to use symbolic variables in matrices we use the symbolic ring \texttt{SR}\\[0pt]
cosnider the following:\\[0pt]
\begin{lstlisting}[language=sage,label= ,caption= ,captionpos=b,numbers=none]
v1 = matrix(SR, 4, 1, [1,2,0,0])
v2 = matrix(SR, 4, 1, [0,1,0,0])
v3 = matrix(SR, 4, 1, [var(f'x{i}') for i in range(4)])
ld(v1.augment(v2).augment(v3), v1.augment(v2).augment(v3).echelon_form())
\end{lstlisting}

\[\left[\begin{array}{rrr} 1 & 0 & x_{0} \\ 2 & 1 & x_{1} \\ 0 & 0 & x_{2} \\ 0 & 0 & x_{3} \end{array}\right] \left[\begin{array}{rrr} 1 & 0 & 0 \\ 0 & 1 & 0 \\ 0 & 0 & 1 \\ 0 & 0 & 0 \end{array}\right]\]\\[0pt]

this isnt really the result we're looking for, because the correct result will most certainly contain the symbolic variables \(x_{0\rightarrow 3}\)\\[0pt]
to get proper results we use \texttt{PolynomialRing}\\[0pt]

\begin{lstlisting}[language=sage,label= ,caption= ,captionpos=b,numbers=none]
R = PolynomialRing(QQ, 4, 'x')
v1 = matrix(R, 4, 1, [1,2,0,0])
v2 = matrix(R, 4, 1, [0,1,0,0])
v3 = matrix(R, 4, 1, R.gens())
ld(v1.augment(v2).augment(v3), v1.augment(v2).augment(v3).echelon_form())
\end{lstlisting}

\[\left[\begin{array}{rrr} 1 & 0 & x_{0} \\ 2 & 1 & x_{1} \\ 0 & 0 & x_{2} \\ 0 & 0 & x_{3} \end{array}\right] \left[\begin{array}{rrr} 1 & 0 & x_{0} \\ 0 & 1 & -2 x_{0} + x_{1} \\ 0 & 0 & x_{2} \\ 0 & 0 & x_{3} \end{array}\right]\]\\[0pt]

now this, this is a proper result, cheers!\\[0pt]
to get the same functionality but with other Rings/Fields we can use \texttt{MatrixSpace}\\[0pt]
here im using the \texttt{SR} field (short for Symbolic Ring)\\[0pt]

\begin{lstlisting}[language=sage,label= ,caption= ,captionpos=b,numbers=none]
R.<a> = PolynomialRing(QQ)
MS = MatrixSpace(SR, 3, 4)
m = MS.matrix([[1,1,1,a],
               [1,a,3,-2],
               [1,3,-1,0]])
fm(r'\[{m} \Longrightarrow {m.echelon_form()} \Longrightarrow {m.echelon_form().simplify_full()}\]')
\end{lstlisting}

\[\left[\begin{array}{rrrr} 1 & 1 & 1 & a \\ 1 & a & 3 & -2 \\ 1 & 3 & -1 & 0 \end{array}\right] \Longrightarrow \left[\begin{array}{rrrr} 1 & 0 & 0 & a + \frac{{\left(a - \frac{2 \, {\left(a + 2\right)}}{a - 1}\right)} {\left(\frac{2}{a - 1} - 1\right)}}{2 \, {\left(\frac{2}{a - 1} + 1\right)}} + \frac{a + 2}{a - 1} \\ 0 & 1 & 0 & -\frac{a + 2}{a - 1} - \frac{a - \frac{2 \, {\left(a + 2\right)}}{a - 1}}{{\left(a - 1\right)} {\left(\frac{2}{a - 1} + 1\right)}} \\ 0 & 0 & 1 & \frac{a - \frac{2 \, {\left(a + 2\right)}}{a - 1}}{2 \, {\left(\frac{2}{a - 1} + 1\right)}} \end{array}\right] \Longrightarrow \left[\begin{array}{rrrr} 1 & 0 & 0 & \frac{1}{2} \, a + 4 \\ 0 & 1 & 0 & -2 \\ 0 & 0 & 1 & \frac{1}{2} \, a - 2 \end{array}\right]\]\\[0pt]

\subsection{vector spaces}
\label{sec:orgcf27a80}
\begin{lstlisting}[language=sage,label= ,caption= ,captionpos=b,numbers=none]
V = VectorSpace(QQ, 4)
S = V.subspace([V([-1,1,2,-2]),V([0,2,-4,2])])
U = V.subspace([V([1,-1,-1,1]),V([2,0,-4,2])])
U.basis()
U.dimension()
S.basis()
S.dimension()
\end{lstlisting}

[\\[0pt]
(1, 0, -2, 1),\\[0pt]
(0, 1, -1, 0)\\[0pt]
]\\[0pt]
2\\[0pt]
[\\[0pt]
(1, 0, -4, 3),\\[0pt]
(0, 1, -2, 1)\\[0pt]
]\\[0pt]
2\\[0pt]

\section{boolean algebra}
\label{sec:org51174d3}
\begin{lstlisting}[language=Python,label= ,caption= ,captionpos=b,numbers=none]
from sympy.logic import SOPform
from sympy import symbols
from sympy.logic.boolalg import truth_table

w, x, y, z = symbols('w x y z')
minterms = [5, 7, 8, 9, 10, 15]
sop = SOPform([w, x, y, z], minterms, [])

return list(truth_table(sop, [w,x,y,z]))
\end{lstlisting}

\begin{tabular}{ll}
(0 0 0 0) & False\\[0pt]
(0 0 0 1) & False\\[0pt]
(0 0 1 0) & False\\[0pt]
(0 0 1 1) & False\\[0pt]
(0 1 0 0) & False\\[0pt]
(0 1 0 1) & True\\[0pt]
(0 1 1 0) & False\\[0pt]
(0 1 1 1) & True\\[0pt]
(1 0 0 0) & True\\[0pt]
(1 0 0 1) & True\\[0pt]
(1 0 1 0) & True\\[0pt]
(1 0 1 1) & False\\[0pt]
(1 1 0 0) & False\\[0pt]
(1 1 0 1) & False\\[0pt]
(1 1 1 0) & False\\[0pt]
(1 1 1 1) & True\\[0pt]
\end{tabular}
\section{latex parsing and integration with emacs}
\label{sec:org6faee3c}
apparently sagemath itself cannot parse latex, but sympy which is packaged with sagemath has an experimental latex parser which should be sufficient, apparently the parser requires a package called \texttt{antlr4-python3-runtime}, i tried installing it globally using pacman but i was still getting the following error\\[0pt]

\begin{lstlisting}[language=Python,label= ,caption= ,captionpos=b,numbers=none]
---------------------------------------------------------------------------
ImportError                               Traceback (most recent call last)
Input In [18], in <cell line: 1>()
----> 1 parse_latex(r"\int_a^b f(x) dx")

File /usr/lib/python3.10/site-packages/sympy/parsing/latex/__init__.py:35, in parse_latex(s)
     30 _latex = import_module(
     31     'sympy.parsing.latex._parse_latex_antlr',
     32     import_kwargs={'fromlist': ['X']})
     34 if _latex is not None:
---> 35     return _latex.parse_latex(s)

File /usr/lib/python3.10/site-packages/sympy/parsing/latex/_parse_latex_antlr.py:65, in parse_latex(sympy)
     62 antlr4 = import_module('antlr4', warn_not_installed=True)
     64 if None in [antlr4, MathErrorListener]:
---> 65     raise ImportError("LaTeX parsing requires the antlr4 Python package,"
     66                       " provided by pip (antlr4-python2-runtime or"
     67                       " antlr4-python3-runtime) or"
     68                       " conda (antlr-python-runtime)")
     70 matherror = MathErrorListener(sympy)
     72 stream = antlr4.InputStream(sympy)

ImportError: LaTeX parsing requires the antlr4 Python package, provided by pip (antlr4-python2-runtime or antlr4-python3-runtime) or conda (antlr-python-runtime)
\end{lstlisting}

so i decided to use \href{https://github.com/OrangeX4/latex2sympy}{latex2sympy}\\[0pt]

install it using:\\[0pt]

\begin{lstlisting}[language=bash,label= ,caption= ,captionpos=b,numbers=none]
pip install latex2sympy2
\end{lstlisting}

example usage:\\[0pt]

\begin{lstlisting}[language=Python,label= ,caption= ,captionpos=b,numbers=none]
from latex2sympy2 import latex2sympy, latex2latex

tex = r"\frac{d}{dx}(x^{2}+x)"
# Or you can use '\mathrm{d}' to replace 'd'
print(latex2sympy(tex))
print(latex2latex(tex))
\end{lstlisting}
\section{more resources}
\label{sec:orga1e193a}
in sage: \url{https://cocalc.com/share/public\_paths/6575abaa81222f77110e4e1de268ff8b12145be8}\\[0pt]
\end{document}